Everything above reduces to a short checklist. It costs one extra column in a
results table and no extra model forwards beyond the ones already being run.

\begin{enumerate}\itemsep3pt
\item \textbf{Write down the hypothesis class before you look at the data, and
report its size.} ``All subsets of these 26 heads'', ``all (layer, position,
$k$-dimensional coordinate block) triples'' --- whatever it is. If you cannot
name it, you cannot correct for it, and you should hold data out instead.

\item \textbf{Say which population your interventions represent.} New fills of
your templates, or new templates? The two give different error bars
(Section~\ref{sec:cluster}); report the cluster-robust one if you mean the
second, and say which you mean.

\item \textbf{Keep the per-instance scores, not just their average.} The
$n\times|\Cclass|$ matrix is what every bound in this paper consumes; it costs
nothing extra to save and it lets a reader recompute any interval.

\item \textbf{Freeze the ablation values.} Mean ablation replaces a component by
an average, and if that average is computed from the evaluation sample itself
then the per-instance score is not a function of a single instance and none of
the bounds apply as stated. Compute the ablation values once, on a reference set
disjoint from the evaluation instances, and hold them fixed. (Resample ablation
has the same issue and is not covered by our guarantees at all.)

\item \textbf{Pick a bound.}
\emph{(i)} If the class was fixed in advance and is enumerable, use the
bootstrap-$t$ simultaneous band --- it is the tightest, it certifies every
hypothesis you looked at, and it adapts to the correlation between them. If some
hypotheses may be perfect on every sampled instance, use the floored variant.
\emph{(ii)} If the class is a power set of components, use the Occam bound with
the size-stratified prior.
\emph{(iii)} If the search was adaptive (greedy pruning, gradient-based alignment
search, or an investigator iterating), hold interventions out. In
interpretability the evaluation data is synthetic and nearly free, so this is
almost always the right answer.

\item \textbf{Report the certified lower bound next to the point estimate}, and
report $m_{\mathrm{eff}}$. A reader then sees both the number and its warranty.

\item \textbf{Pre-register the final circuit.} The cheapest way to get a tight
interval is to fix the circuit, generate fresh interventions, and evaluate once:
the multiplicity penalty disappears entirely. In our GPT-2 experiment this is
worth \MyPreRegGain{} of certified faithfulness at $n=1000$
(Table~\ref{tab:main} versus the pre-registered row in
Appendix~\ref{app:extra}).
\end{enumerate}
